% ============================================================================
% Covariant-Complete Dual Tetrad Gravity
% Supplemental Material to "Dual Tetrad Gravity And The Field Origin"
% ============================================================================
\documentclass[11pt]{article}

\usepackage{amsmath,amssymb}
\usepackage{bm}
\usepackage{geometry}
\geometry{margin=1in}
\usepackage{hyperref}
\usepackage{booktabs}
\usepackage{xcolor}
\usepackage{enumitem}

\hypersetup{
    colorlinks=true,
    linkcolor=blue!60!black,
    citecolor=blue!60!black,
    urlcolor=blue!60!black
}

\setlength{\parindent}{0pt}
\setlength{\parskip}{0.8ex plus 0.2ex minus 0.1ex}

% ============================================================================
\begin{document}

\begin{center}
{\Large\bfseries Covariant-Complete Dual Tetrad Gravity:\\[4pt] A Zero-Parameter Lagrangian from Topological Field Closure}\\[12pt]
{\normalsize Stephen Nelson}\\
{\small Independent Researcher}\\[8pt]
{\small February 2026}\\[6pt]
{\small\itshape Supplemental Material to ``Dual Tetrad Gravity And The Field Origin''}
\end{center}

\vspace{0.5cm}
\hrule
\vspace{0.5cm}

\textbf{Summary.}\quad
A central open question in gravitational physics is whether the coupling constants of gravity are free parameters or whether they follow from structural requirements of the field theory itself.
The dual tetrad coupling topology answers this question affirmatively.
Two tetrahedra with opposite orientations coupled at a shared center form the minimal structure that achieves full field closure in three spatial dimensions.
This topology forces a finite coupling channel count, which imposes a saturation condition on the gravitational field; saturation, combined with detailed balance, uniquely determines the constitutive law $\mu(x) = x/(1+x)$; and the constitutive law integrates to a unique scalar Lagrangian through the standard AQUAL variational framework~[6].
The result is a covariant-complete, zero-parameter scalar-tensor action whose field equations produce five quantitative predictions across gravity, cosmology, particle physics, galactic dynamics, and the CMB.
This supplemental note presents the action, its derivation from the topology, its field equations, and its predictions in self-contained form, accompanying the manuscript ``Dual Tetrad Gravity And The Field Origin''~[7].

% ============================================================================
\section*{I.\quad The Dual Tetrad Coupling Structure}
% ============================================================================

Tetrad gravity describes the gravitational field through local frames (vierbeins) whose spatial triad defines a tetrahedron, the minimal solid in three dimensions.
Ghost-free bimetric gravity couples two such tetrad fields through a frame transition matrix $\mathcal{S}^a_{\;b}$ and an interaction potential with five parameters $\beta_0$ through $\beta_4$~[5].
The manuscript~[7] proves that this frame transition matrix is identically the coupled face matrix of the dual tetrad topology.
The identification rests on three results: the triad-tetrahedron correspondence, the symmetric vielbein condition, and the compatible $(3{+}1)$ decompositions required for general covariance.
Before presenting the action, we summarize the structural origin of the coupling topology, because the entire Lagrangian follows from it.

The ADM decomposition of General Relativity imposes $k = 4$ constraint equations at each spatial point: one Hamiltonian constraint plus three momentum constraints~[1].
Each constraint governs the gravitational coupling across one interface of a field region.
The minimum number of such interfaces in $d$ spatial dimensions is $d + 1$, the face count of the simplex: in $d = 3$, a tetrahedron has $3 + 1 = 4$ faces.
The ADM constraint count and the simplex face count agree because both answer the same question: how many independent coupling directions does a gravitational field region have?

Why $d = 3$ and not some other dimension?
Two tetrahedra coupled at a shared center must account for every independent direction in the embedding space; otherwise, some coupling directions are unresolved and the field cannot close.
In $d$ dimensions, two simplices contribute $2(d+1)$ vertices, and the embedding hypercube has $2^d$ corners.
Closure requires these to match: $2(d+1) = 2^d$.
This equation has a unique integer solution: $d = 3$, giving $2 \times 4 = 8 = 2^3$.
Three spatial dimensions are the unique case where dual simplicial field closure accounts for every direction in the embedding space.

Having established $k = 4$, the interaction count follows.
Two tetrahedra with opposite orientations sharing a common center form the stellated octahedron~[2].
In this structure, the tetrahedra interpenetrate: each of the 4 vertices of one tetrahedron passes through a face of the other.
With 4 vertices penetrating 4 faces, there are $k \times k = (3{+}1)^2 = 16$ independent coupling channels.
This dual tetrad topology is the structure from which the Lagrangian, the acceleration scale, and the gravitational coupling are derived.

\textbf{The gravitational coupling.}\quad
The counting above determines $G$ directly:
\begin{equation}
G = \frac{(3+1)^2 + 1}{3(3+1) + 1} \times \alpha^{\,1 + (3+1) + (3+1)^2} \times \frac{\hbar c}{m_e^2} \;=\; \frac{17}{13}\,\alpha^{21}\,\frac{\hbar c}{m_e^2}
\label{eq:G}
\end{equation}
The exponent $21 = 1 + k + k^2 = 1 + 4 + 16$ counts the total coupling depth: 1 for the interaction vertex (Field Origin) where the two tetrahedra meet, $k = 4$ for the coupling interfaces of one sector, and $k^2 = 16$ for the full bimetric interaction.
The prefactor $17/13$ is the ratio of active coupling modes ($16$ channels $+ 1$ vertex $= 17$) to structural elements ($12$ edges $+ 1$ center $= 13$).
Every integer is counted from the stellated octahedron; none is adjustable.
The full counting derivation, with a term-by-term translation table mapping each integer to its GR origin, is presented in the companion note ``The Origin of Every Integer in $G$.''

\textbf{Field closure and the dual tetrad topology are the same structure.}\quad
Field closure is the physical requirement that a distributed gravitational field settle to a consistent configuration at every point.
In the language of GR, it is what the ADM constraint equations enforce: energy and momentum must balance across every coupling interface~[1].
The dual tetrad is the topological answer to the question: what is the minimum structure that achieves full field closure in $d = 3$ spatial dimensions?
The answer is two $k = 4$ simplices coupled at a shared center, producing $k^2 = 16$ coupling channels.
The topology does not merely ``require'' field closure; the topology \emph{is} field closure, the unique form that the consistency requirement takes in three-dimensional space.

This identification is what connects the Lagrangian (derived in Section~III) to the topology (established above).
The dual tetrad topology forces field closure.
Field closure, acting on a finite channel count, forces the constitutive law $\mu(x) = x/(1+x)$ through saturation and detailed balance.
The constitutive law integrates to a unique scalar Lagrangian $\mathcal{F}(y)$.
The topology, the physical requirement, and the action are three descriptions of one structure.

% ============================================================================
\section*{II.\quad The Overdetermined Consistency Test}
% ============================================================================

Having established where the integers originate, we now ask: does the resulting framework capture real physical structure?
The most stringent test available is overdetermination.
When more constraints exist than unknowns, a solution exists \emph{only if} the constraints are mutually consistent.
This is why GPS receivers use four satellites to fix three coordinates; the fourth is not redundant but a consistency check.
The dual tetrad framework passes this test.

Three equations from three independent domains constrain two unknowns ($r$ and $m$, the size and mass of a single tetrahedral field region):

\begin{enumerate}[leftmargin=2em]
\item \textbf{Gravitational closure:}\quad $a_0 = k^2\,Gm/r^2$\quad (where $a_0$ is the acceleration scale below which gravitational dynamics deviate from Newtonian behavior)
\item \textbf{Cosmological horizon:}\quad $a_0 = cH/(2\pi)\sqrt{k/\pi}$\quad (from thermodynamic arguments applied to the Hubble horizon)
\item \textbf{Quantum-electromagnetic scale:}\quad $m = \alpha\hbar/(rc)$\quad (the classical electron radius relation)
\end{enumerate}

Three equations, two unknowns.
Most overdetermined systems have no solution at all.

The system converges:
\begin{equation*}
r = 2.82 \times 10^{-15}\text{ m}\,,\qquad m = 9.11 \times 10^{-31}\text{ kg}
\end{equation*}
matching the measured electron radius ($r_e = 2.818$ fm) and electron mass ($m_e = 9.109 \times 10^{-31}$ kg) to 0.1\%, with zero adjustable parameters.
The inputs are: $k = 4$ (simplex theorem), $G$ (Eq.~\ref{eq:G}, derived from the topology to 0.12\%), $\alpha$ (the fine-structure constant; its role in gravity is derived in the companion note ``Why the Fine-Structure Constant Governs Gravitational Coupling''), $\hbar$, $c$, and $H$ (a cosmological observable independent of particle physics).
No particle-physics data enters as input; the electron is an \emph{output}.

\textbf{Bidirectional closure.}\quad
Starting from $m_e$, the framework derives $H = 69.92$ km/s/Mpc, matching the gravitational-wave standard siren measurement ($70.0 \pm 12$)~[3] to 0.11\%.
Reversing the calculation, the observed $H$ recovers $m_e$ to 0.10\%.
This round-trip consistency across particle and cosmological scales demonstrates structural correctness: a unidirectional fit cannot close bidirectionally with comparable precision.

% ============================================================================
\section*{III.\quad The Scalar-Tensor Action from Topological Field Closure}
% ============================================================================

Because the dual tetrad topology \emph{is} field closure (Section~I), it directly determines a scalar-tensor action through the following chain.
The finite channel count ($k^2 = 16$) imposes a saturation condition on gravitational coupling: at weak field gradients, the field cannot recruit channels beyond those available, and the response flattens.
Saturation, combined with detailed balance, determines a unique constitutive law.
The constitutive law then integrates to a unique Lagrangian via the AQUAL variational relationship~[6].
The result is a covariant action with disformal electromagnetic coupling in which every component (the Lagrangian, the acceleration scale, and the gravitational coupling) is derived from the topology.
The action contains zero free parameters.

\begin{equation}
S = \int d^4x \sqrt{-g} \left[ \frac{R}{16\pi G} - \frac{a_0^2}{8\pi G}\,\mathcal{F}\!\left(\frac{|\nabla\Phi|^2}{a_0^2}\right) - \frac{1}{4}\,e^{-2\Phi/c^2} F_{\mu\nu}F^{\mu\nu} + \mathcal{L}_{\text{matter}} \right]
\label{eq:action}
\end{equation}

Each term has a definite origin:
\begin{itemize}[leftmargin=2em]
\item $R/(16\pi G)$: spacetime curvature (Einstein--Hilbert), with $G$ given by Eq.~\ref{eq:G}
\item $\mathcal{F}(y)$: gravitational saturation from field closure (derived below)
\item $e^{-2\Phi/c^2} F_{\mu\nu}F^{\mu\nu}$: disformal electromagnetic coupling; photons propagate on the effective metric $\tilde{g}_{\mu\nu} = e^{-2\Phi/c^2}\,g_{\mu\nu}$, recovering the full GR deflection angle $\theta = 4GM/(rc^2)$ in the weak-field limit
\item $\mathcal{L}_{\text{matter}}$: baryonic sources
\end{itemize}

\textbf{From topology to Lagrangian: the derivation chain.}\quad
As established in Section~I, the dual tetrad topology is the structural form of field closure: gravity's requirement to define a complete, self-consistent configuration at every point.
In the dual tetrad structure, all 16 coupling channels must be accounted for; no coupling direction may be left unresolved.

Field closure imposes a saturation condition on the gravitational field.
At strong field gradients (well above the topological scale $a_0$), all coupling channels are fully active and the field responds linearly; this is the Newtonian regime.
At weak field gradients (well below $a_0$), the finite number of coupling channels saturates: the field cannot recruit additional channels, and the gravitational response flattens.
The requirement that the transition between these regimes satisfy detailed balance (the rate of channel activation equals the rate of saturation) uniquely determines the constitutive law:
\begin{equation*}
\mu(x) = \frac{x}{1+x}
\end{equation*}
where $x = |\nabla\Phi|/a_0$ measures the field gradient in units of the topological scale.
This function has the correct limits: $\mu \to 1$ (Newtonian) when $x \gg 1$, and $\mu \to x$ (saturated) when $x \ll 1$.

The scalar Lagrangian $\mathcal{F}$ is then \emph{derived} by integrating this constitutive law.
The Euler--Lagrange equation for the scalar field $\Phi$ relates the Lagrangian to the constitutive law through the condition $d\mathcal{F}/dx = x\,\mu(x)$ (the standard AQUAL variational relationship~[6]), where $y = x^2 = |\nabla\Phi|^2/a_0^2$.
Integrating:
\begin{equation}
\frac{d\mathcal{F}}{dx} = \frac{x^2}{1+x} \quad\longrightarrow\quad \mathcal{F}(y) = \tfrac{1}{2}y - \sqrt{y} + \ln(1 + \sqrt{y})
\label{eq:lagrangian}
\end{equation}
This integration is unique: the integration constant is fixed by $\mathcal{F}(0) = 0$, and the constitutive law admits no alternative form under the field closure boundary condition.
The Lagrangian is the unique output of a requirement that the topology imposes.
The full derivation of $\mu(x)$ from the stellated octahedron's coupling geometry appears in the companion note ``Why the Fine-Structure Constant Governs Gravitational Coupling,'' Section~V.

\textbf{Limit checks.}\quad
In the Newtonian regime ($y \gg 1$): $\mathcal{F}(y) \to \tfrac{1}{2}y$, recovering the standard kinetic term and reproducing GR.
In the deep-MOND regime ($y \ll 1$): $\mathcal{F}(y) \to \sqrt{y}$, producing the modified dynamics observed in galactic rotation curves~[4].
The logarithmic term encodes the smooth crossover between regimes.
The action reproduces GR where GR is well-tested and produces MOND phenomenology where dark matter is otherwise required, with correct asymptotic behavior in both limits.

\textbf{Field equations.}\quad
Varying the action (Eq.~\ref{eq:action}) with respect to the scalar field and the metric produces three sets of field equations:
\begin{enumerate}[leftmargin=2em]
\item A \emph{modified Poisson equation} governing galactic-scale gravitational dynamics
\item \emph{Modified Friedmann equations} governing cosmological expansion
\item A \emph{growth equation} governing cosmological structure formation
\end{enumerate}
These field equations predict the CMB gravitational enhancement factor (the ratio of total gravitating mass to baryonic mass at recombination) as $1/\mu(x_{\text{rec}}) = 6.46$, matching Planck 2018's measured value of 6.40 to 1\%, with zero free parameters.

\textbf{There are zero free parameters in this action.}
The Lagrangian $\mathcal{F}$ (Eq.~\ref{eq:lagrangian}), the acceleration scale $a_0 = k^2 G m_e / r_e^2$, and the gravitational coupling $G$ are all determined by the dual tetrad structure.
Unlike $\Lambda$CDM, which requires six fitted cosmological parameters and a dark matter component, this action is fully specified by the topology.

% ============================================================================
\section*{IV.\quad Zero-Parameter Predictions}
% ============================================================================

The action and its field equations produce the following predictions:

\begin{center}
\begin{tabular}{lp{4.2cm}p{3.5cm}c}
\toprule
\textbf{Quantity} & \textbf{Predicted} & \textbf{Measured} & \textbf{Deviation} \\
\midrule
$G$ (gravitational constant) & $6.666 \times 10^{-11}$\newline m$^3$kg$^{-1}$s$^{-2}$ & $6.674 \times 10^{-11}$ & 0.12\% \\[4pt]
$H$ (Hubble constant) & 69.92 km/s/Mpc & $70.0 \pm 12$ & 0.11\% \\[4pt]
$m_e$ (electron mass) & $9.11 \times 10^{-31}$ kg & $9.109 \times 10^{-31}$ & 0.10\% \\[4pt]
$a_0$ (MOND scale) & $1.20 \times 10^{-10}$ m/s$^2$ & $1.20 \times 10^{-10}$ & 0.41\% \\[4pt]
CMB enhancement & 6.46 & 6.40 (Planck 2018) & 1.0\% \\
\bottomrule
\end{tabular}
\end{center}
\vspace{0.3em}

\noindent
Five predictions across five domains, all with zero adjustable parameters.
Changing the coupling exponent from 21 to 20 or 22 shifts the predicted $G$ by a factor of 137 ($= 1/\alpha$); only the counted value produces agreement.

% ============================================================================
\section*{V.\quad Structural Consistency Tests}
% ============================================================================

The framework exhibits the following structural properties, each independently testable:

\begin{enumerate}[leftmargin=2em]
\item \textbf{Overdetermined convergence:} Three independent equations from three domains converge on a single solution (Section~II). Overdetermined systems generically have no solution; convergence indicates mutual consistency of the constraints.

\item \textbf{Bidirectional closure:} Deriving $H$ from $m_e$ and recovering $m_e$ from $H$, each to $\sim$0.1\%, with the same topology. A unidirectional fit cannot close bidirectionally with comparable precision.

\item \textbf{Sharp selectivity:} Changing the exponent $21$ by $\pm 1$ shifts $G$ by exactly $1/\alpha = 137.036\ldots$, because $G \propto \alpha^n$. The topology determines both the exponent and the selectivity. Only the counted value produces agreement. (The full sensitivity analysis appears in the companion note, ``The Origin of Every Integer in $G$.'')

\item \textbf{Correct asymptotic limits:} The Lagrangian $\mathcal{F}(y)$ reduces to the standard kinetic term ($\tfrac{1}{2}y$) in the Newtonian regime and to the MOND Lagrangian ($\sqrt{y}$) in the deep-field regime. These are structural properties of the derived Lagrangian.

\item \textbf{Determined, falsifiable action:} The action (Eq.~\ref{eq:action}) is fully specified. Its field equations yield cosmological predictions (CMB enhancement to 1\%). The theory can be falsified by solving these field equations and comparing to observations. This is the standard of a covariant field theory.

\item \textbf{Invertibility:} The constitutive law $\mu(x) = x/(1+x)$ is bijective: given dynamics, one can uniquely recover baryonic mass, and vice versa. In $\Lambda$CDM, the reverse map (dynamics $\to$ baryonic mass) requires assuming a dark matter halo profile with additional free parameters. This structural property is testable galaxy by galaxy.
\end{enumerate}

% ============================================================================
\section*{VI.\quad Acknowledged Limitations}
% ============================================================================

The framework does not address:
\begin{itemize}[leftmargin=2em]
\item Galaxy cluster dynamics (the Bullet Cluster and residual cluster mass)
\item The full CMB angular power spectrum shape (the enhancement factor is predicted to 1\%, but the detailed multipole structure has not been computed)
\item The strong and weak nuclear forces
\item The cosmological constant
\end{itemize}

\noindent
The 0.12\% residual in $G$ exceeds the $\sim$0.05\% inter-experiment scatter, suggesting either higher-order corrections or systematic measurement effects.
The manuscript~[7] also demonstrates that the dual tetrad topology constrains the coupling parameters of the established Hassan--Rosen bimetric action~[5]; the relationship between that bimetric formulation and the scalar-tensor action presented here is identified as a target for future work.

These limits are stated in the manuscript~[7].
The goal is clarity about what the topology explains and where the open questions lie.

% ============================================================================
\vspace{0.8cm}
\hrule
\vspace{0.5cm}

\small
\textbf{References}\\[4pt]
{[1]} R.~Arnowitt, S.~Deser, and C.~W.~Misner, ``The Dynamics of General Relativity,'' Gen.\ Relativ.\ Gravit.\ \textbf{40}, 1997 (2008).\\
{[2]} H.~S.~M.~Coxeter, \textit{Regular Polytopes}, 3rd ed.\ (Dover, 1973).\\
{[3]} B.~P.~Abbott \textit{et al.} (LIGO/Virgo), ``A gravitational-wave standard siren measurement of the Hubble constant,'' Nature \textbf{551}, 85 (2017).\\
{[4]} M.~Milgrom, ``A modification of the Newtonian dynamics as a possible alternative to the hidden mass hypothesis,'' Astrophys.\ J.\ \textbf{270}, 365 (1983).\\
{[5]} S.~F.~Hassan and R.~A.~Rosen, ``Bimetric gravity from ghost-free massive gravity,'' JHEP \textbf{02}, 126 (2012).\\
{[6]} J.~Bekenstein and M.~Milgrom, ``Does the missing mass problem signal the breakdown of Newtonian gravity?'' Astrophys.\ J.\ \textbf{286}, 7 (1984).\\
{[7]} S.~Nelson, ``Dual Tetrad Gravity And The Field Origin'' (2026).
Submitted for publication.

\end{document}
